% Generado a partir de calculus/Final_Bien/evidencia_P8.m (simulador corregido).
\subsection{Sistema no lineal en lazo abierto en los puntos de equilibrio}

Para caracterizar la inestabilidad de la configuración atractiva, se simuló el modelo no lineal sin fricción estática ni zona muerta, aplicando el voltaje de equilibrio $u=u_e(x_1^*)$ y desplazando la posición inicial $0{,}01$~cm respecto al equilibrio, en $x_1^*=-4$, $-3$ y $-2$~cm. La figura~\ref{fig:p8_lazo_abierto} muestra que la desviación crece exponencialmente y alcanza 1~cm en $0{,}33$, $0{,}32$ y $0{,}30$~s, respectivamente; el modelo linealizado reproduce el crecimiento con un tiempo de duplicación $\ln 2/\lambda_2$ de $46$, $44$ y $41$~ms. Cuanto más cerca de la bobina se encuentra el equilibrio, más rápido es el polo inestable.

\begin{figure}[ht]
    \centering
    % Estilo común de las figuras de evidencia P8 (pgfplots).
\pgfplotsset{
  pocho/.style={
    width=0.92\textwidth, height=5.2cm,
    grid=major, grid style={gray!25},
    tick label style={font=\small}, label style={font=\small}, legend style={font=\small},
    /pgf/number format/use comma,
    scaled ticks=false,
    y tick label style={/pgf/number format/fixed},
    table/col sep=comma, unbounded coords=jump,
    every axis plot/.append style={line join=round},
  },
}
\definecolor{tray}{RGB}{31,94,166}
\definecolor{refc}{RGB}{40,40,40}
\definecolor{ctl2}{RGB}{196,96,40}
\definecolor{ver}{RGB}{46,133,64}

\begin{tikzpicture}
\begin{semilogyaxis}[pocho, height=6cm, xmin=0, xmax=0.36, ymin=0.01, ymax=1,
    xlabel={tiempo [\si{s}]}, ylabel={$|x_1(t)-x_1^*|$ [\si{cm}]},
    x tick label style={/pgf/number format/fixed},
    legend style={at={(0.02,0.98)}, anchor=north west}, legend cell align=left]
  \addplot[tray, line width=1pt]  table[x=t, y=nl_m4] {images/p8_results/tikz/lazo_abierto.csv}; \addlegendentry{$x_1^*=-4$ cm}
  \addplot[ctl2, line width=1pt]  table[x=t, y=nl_m3] {images/p8_results/tikz/lazo_abierto.csv}; \addlegendentry{$x_1^*=-3$ cm}
  \addplot[ver, line width=1pt]   table[x=t, y=nl_m2] {images/p8_results/tikz/lazo_abierto.csv}; \addlegendentry{$x_1^*=-2$ cm}
  \addplot[tray, densely dashed]  table[x=t, y=lin_m4] {images/p8_results/tikz/lazo_abierto.csv}; \addlegendentry{modelo lineal}
  \addplot[ctl2, densely dashed, forget plot] table[x=t, y=lin_m3] {images/p8_results/tikz/lazo_abierto.csv};
  \addplot[ver, densely dashed, forget plot]  table[x=t, y=lin_m2] {images/p8_results/tikz/lazo_abierto.csv};
\end{semilogyaxis}
\end{tikzpicture}

    \caption[Lazo abierto en los puntos de equilibrio]{Desviación respecto al equilibrio en lazo abierto, con $u=u_e$ y una perturbación inicial de $0{,}01$~cm. Líneas continuas: modelo no lineal; discontinuas: modelo linealizado.}
    \label{fig:p8_lazo_abierto}
\end{figure}

El retrato de fase alrededor de $x_1^*=-4$~cm (figura~\ref{fig:p8_retrato}) confirma que el equilibrio es un punto silla: las trayectorias se alejan a lo largo de la variedad inestable, salvo las que parten exactamente sobre la variedad estable. Al incorporar la fricción estática de Karnopp cambia la situación en reposo: si el levitante está detenido y $u=u_e$, la fuerza neta $u/[b(a-x)^4]-mg$ no supera $F_s$ mientras la posición se encuentre entre $-4{,}44$ y $-3{,}63$~cm, de modo que el imán permanece atascado en cualquier punto de esa banda. Para $x_1^*=-3$ y $-2$~cm las bandas correspondientes son $[-3{,}41;\,-2{,}66]$ y $[-2{,}37;\,-1{,}70]$~cm. La fricción estática enmascara así la inestabilidad del equilibrio mientras el levitante no se mueva.

\begin{figure}[ht]
    \centering
    % Estilo común de las figuras de evidencia P8 (pgfplots).
\pgfplotsset{
  pocho/.style={
    width=0.92\textwidth, height=5.2cm,
    grid=major, grid style={gray!25},
    tick label style={font=\small}, label style={font=\small}, legend style={font=\small},
    /pgf/number format/use comma,
    scaled ticks=false,
    y tick label style={/pgf/number format/fixed},
    table/col sep=comma, unbounded coords=jump,
    every axis plot/.append style={line join=round},
  },
}
\definecolor{tray}{RGB}{31,94,166}
\definecolor{refc}{RGB}{40,40,40}
\definecolor{ctl2}{RGB}{196,96,40}
\definecolor{ver}{RGB}{46,133,64}

\begin{tikzpicture}
\begin{axis}[pocho, width=0.8\textwidth, height=7.5cm, xmin=-4.5, xmax=-3.5, ymin=-10, ymax=10,
    xlabel={posición $x_1$ [\si{cm}]}, ylabel={velocidad $x_2$ [\si{cm/s}]},
    legend style={at={(0.98,0.02)}, anchor=south east, fill=white}, legend cell align=left]
  \addplot[gray!60, quiver={u=\thisrow{dx}, v=\thisrow{dv}, every arrow/.append style={-stealth}}, forget plot]
    table[x=x, y=v] {images/p8_results/tikz/campo.csv};
  \addplot[tray, line width=0.9pt] table[x=x, y=v] {images/p8_results/tikz/trayectorias.csv}; \addlegendentry{trayectorias}
  \addplot[ctl2, line width=1.4pt] table[x=x, y=v] {images/p8_results/tikz/variedad_inestable.csv}; \addlegendentry{variedad inestable}
  \addplot[ver, line width=3pt, opacity=0.6] coordinates {(-4.4448,0) (-3.6291,0)}; \addlegendentry{banda de atascamiento ($u=u_e$)}
  \addplot[only marks, mark=*, mark size=2pt, refc] coordinates {(-4,0)}; \addlegendentry{equilibrio $x_1^*=-4$ cm}
\end{axis}
\end{tikzpicture}

    \caption[Retrato de fase en lazo abierto]{Retrato de fase del modelo no lineal sin fricción estática alrededor de $x_1^*=-4$~cm, con $u=u_e$. La banda verde indica las posiciones en las que, con el levitante en reposo, la fricción estática lo mantiene atascado.}
    \label{fig:p8_retrato}
\end{figure}

\FloatBarrier
